\documentclass[12pt,a4paper]{article}
\usepackage[utf8]{inputenc}
\usepackage[T1]{fontenc}
\usepackage[ngerman]{babel}
\usepackage[margin=2cm,top=3.5cm,bottom=2.5cm]{geometry}
\usepackage{enumitem}
\usepackage{setspace}
\usepackage{parskip}
\usepackage{fancyhdr}
\usepackage{ragged2e}
\usepackage{tikz}
\usepackage{eurosym}

\title{}
\author{}
\date{}

\setlength{\parindent}{0pt}
\setlength{\parskip}{6pt}
\renewcommand{\baselinestretch}{1.15}

\newcommand\Committee{SEDE}
\newcommand\Fraktion{LEER}

\newcommand\Thema{Armee}
\newcommand\city{Ulm}
\newcommand\datum{11.05.2026}
\newcommand\timeVorb{07:30-09:00}
\newcommand\timeEinf{09:00-09:30}
\newcommand\timeFrakOne{09:30-11:00}
\newcommand\timePauseOne{11:00-11:15}
\newcommand\timeAuss{11:15-12:30}
\newcommand\timeMittag{12:30-13:00}
\newcommand\timeFrakTwo{13:00-13:30}
\newcommand\timeFinVerh{13:30-14:00}
\newcommand\timePlenar{14:00-15:15}
\newcommand\timeDebr{15:15-15:30}
\newcommand\politiker{Andrea Wechsler}
\newcommand\politikerOffice{Mitglied des Europäischen Parlaments}
\newcommand\stadtvertreter{Nalan Schmidt}
\newcommand\stadtvertreterOffice{Anlauf- und Koordinierungsstelle Bürgerdialog (Open Government) der Stadt Ulm}
\newcommand\localSupport{das Europe Direct Ulm}
\newcommand\sponsor{die Vertretung der Europäischen Kommission in München}
\newcommand\jefvorsitz{Farras Fathi}
\newcommand\gendervorsitz{Landesvorsitzender}
\newcommand\evpLeader{Tobias}
\newcommand\evpRoom{Kleiner Sitzungssaal}
\newcommand\sdLeader{Antonia}
\newcommand\sdRoom{Veranstaltungsraum Europe Direct}
\newcommand\reLeader{Alica}
\newcommand\reRoom{Besprechungszimmer Donaubüro I}
\newcommand\fifthLeader{Farras}
\newcommand\fifthRoom{Besprechungszimmer Öffentlichkeitsarbeit}
\newcommand\pfeLeader{Hannah}
\newcommand\pfeRoom{Besprechungszimmer BM III}
\newcommand\numSuS{27}
\newcommand\location{im Ulmer Rathaus}
\newcommand\fifthGroup{Linke}
\newcommand\anzahlcomm{3}
\newcommand\totcoms{26}
\newcommand\minmems{25}
\newcommand\maxmems{90}

\newcommand{\LogoBasePath}{../Logos}
% Shared logo helpers for SimEP LaTeX documents.

\providecommand{\LogoBasePath}{Logos}
\providecommand{\FraktionsLogoBasePath}{Folien/Bilder}

\providecommand{\constrainedlogo}[4][2cm]{%
  \raisebox{#3}{\includegraphics[width=#2,height=#1,keepaspectratio]{#4}}%
}

\providecommand{\jeflogo}{\constrainedlogo[2.5cm]{5cm}{0pt}{\LogoBasePath/jef.png}}
\providecommand{\simeplogo}{\constrainedlogo[2.5cm]{5cm}{0pt}{\LogoBasePath/simep.png}}
\providecommand{\financelogo}{\constrainedlogo[1.5cm]{4cm}{0pt}{\LogoBasePath/finance\_\city.png}}
\providecommand{\citylogo}{\constrainedlogo[1.5cm]{4cm}{0pt}{\LogoBasePath/\city.png}}
\providecommand{\fraktionslogo}{\constrainedlogo[2.5cm]{5cm}{0pt}{\FraktionsLogoBasePath/\Fraktion.png}}

\providecommand{\PlaceLogosOne}{%
\begin{tikzpicture}[remember picture, overlay]
        % Top left corner
        \node at ([xshift=10mm, yshift=-8mm] current page.north west) [anchor=north west, inner sep=0pt] {
            {\jeflogo}
        };

        % Top right corner
        \node at ([xshift=-7.5mm, yshift=-6mm] current page.north east) [anchor=north east, inner sep=0pt] {
            {\simeplogo}
        };


        % Bottom left corner
        \node[anchor=south west, inner sep=0pt, text width=4.5cm, align=left] at ([xshift=10mm, yshift=8mm] current page.south west) {
            \small Auftraggeber \& \newline Kooperationspartner:\\[3mm]
            \financelogo
        };


        % Bottom right corner
        \node[anchor=south east, inner sep=0pt, text width=4.5cm, align=right] at ([xshift=-10mm, yshift=10mm] current page.south east) {
            \small Unterst\"utzt durch:\\[3mm]
            \citylogo
        };

\end{tikzpicture}
}

\providecommand{\PlaceLogosRest}{%
\begin{tikzpicture}[remember picture, overlay]
        % Top left corner
        \node at ([xshift=10mm, yshift=-8mm] current page.north west) [anchor=north west, inner sep=0pt] {
            {\jeflogo}
        };

        % Top right corner
        \node at ([xshift=-7.5mm, yshift=-6mm] current page.north east) [anchor=north east, inner sep=0pt] {
            {\simeplogo}
        };

\end{tikzpicture}
}

\providecommand{\PlaceLogosFraktion}{%
\begin{tikzpicture}[remember picture, overlay]
        % Top left corner
        \node at ([xshift=10mm, yshift=-8mm] current page.north west) [anchor=north west, inner sep=0pt] {
            {\jeflogo}
        };

        % Top right corner
        \node at ([xshift=-7.5mm, yshift=-6mm] current page.north east) [anchor=north east, inner sep=0pt] {
            {\simeplogo}
        };

        % Upper-right caucus logo
        \node at ([xshift=-12mm, yshift=-35mm] current page.north east) [anchor=north east, inner sep=0pt] {
            {\fraktionslogo}
        };

        % Bottom left corner
        \node[anchor=south west, inner sep=0pt, text width=4.5cm, align=left] at ([xshift=10mm, yshift=8mm] current page.south west) {
            \small Auftraggeber \& \newline Kooperationspartner:\\[3mm]
            \financelogo
        };


        % Bottom right corner
        \node[anchor=south east, inner sep=0pt, text width=4.5cm, align=right] at ([xshift=-10mm, yshift=10mm] current page.south east) {
            \small Unterst\"utzt durch:\\[3mm]
            \citylogo
        };

\end{tikzpicture}
}

\begin{document}


\thispagestyle{empty}

\vspace*{2cm}

\begin{center}
\Large
Vorschlag für eine

\vspace{0.8cm}

\textbf{\large Verordnung des Europäischen Parlaments und des Rates\\[0.3em]
für ein solidarisches und einheitliches Asylsystem}

\vspace{2cm}

\normalsize
Federführende Ausschüsse:

\vspace{0.5cm}

\textbf{Haushaltsausschuss (BUDG)}

\textbf{Ausschuss für Bürgerliche Freiheiten, Justiz und Inneres (LIBE)}

\textbf{Ausschuss für Beschäftigung und soziale Angelegenheiten (EMPL)}

\textbf{Unterausschuss für Menschenrechte (DROI)}

\end{center}

\PlaceLogosOne


\newpage

\subsection*{IN ERWÄGUNG NACHSTEHENDER GRÜNDE: (Alle Ausschüsse)}

\begin{enumerate}[label=(\arabic*), leftmargin=2em, itemsep=8pt]
    \item Die umfassende Anwendung der Genfer Flüchtlingskonvention vom 28. Juli 1951 stellt die Grundlage eines gemeinsamen Europäischen Asylsystems dar.

    \item Ziel der Europäischen Union ist es, schrittweise einen Raum der Freiheit, der Sicherheit und des Rechts aufzubauen, der allen offensteht, die wegen besonderer Umstände rechtmäßig in der Union Schutz suchen. Ein gemeinsames Europäisches Asylsystem ist hierzu ein wesentlicher Bestandteil.

    \item Für eine solche Politik sollte der Grundsatz der Solidarität und der gerechten Aufteilung der Verantwortlichkeiten unter den Mitgliedstaaten der EU (im Folgenden auch Mitgliedstaaten), auch in finanzieller Hinsicht, gelten.

    \item Erst durch eine gerechte Aufteilung von Lasten und Verantwortlichkeiten können gemeinsame Standards für Asylverfahren, sowie für die Aufnahmebedingungen Asylsuchender gewährleistet und damit auch Sekundärmigration zwischen den Mitgliedstaaten der EU verhindert werden.
\end{enumerate}

\vspace{1cm}

\noindent\textbf{HABEN FOLGENDE VERORDNUNG ERLASSEN:}

\subsection*{Artikel 1 – Gegenstand und Ziele}

\begin{enumerate}[leftmargin=2em, itemsep=6pt]
  \item Die Mitgliedstaaten verpflichten sich zur Realisierung eines solidarischen und einheitlichen Europäischen Asylsystems. Die EU unterstützt die Mitgliedstaaten bei der Umsetzung der geltenden Bestimmungen des gemeinsamen Europäischen Asylsystems. (DROI)
  \item Ein gemeinsamer Asyl- und Integrationsfonds wird eingerichtet, um die finanziellen Lasten im Rahmen der Asyl- und Integrationspolitik gerecht zu verteilen. (BUDG)
  \item Die Europäische Agentur für die Grenz- und Küstenwache (FRONTEX) unterstützt die Mitgliedstaaten bei der Sicherung der Außengrenzen. (LIBE)
  \item Die Mitgliedstaaten fördern die soziale und wirtschaftliche Integration Schutzsuchender durch gezielte Maßnahmen im Bereich Beschäftigung, Bildung und soziale Inklusion. (EMPL)
\end{enumerate}

\subsection*{Artikel 2 – Geltungsbereich}

Die vorliegende Verordnung gilt für alle Maßnahmen der Asyl- und Integrationspolitik, die im Hoheitsgebiet der Mitgliedstaaten der EU realisiert werden. Dies schließt Maßnahmen an den Außengrenzen und in den Hoheitsgewässern mit ein.

\newpage

\subsection*{Artikel 3 – Asyl- und Integrationsfonds (BUDG)}

\begin{enumerate}[leftmargin=2em, itemsep=6pt]
  \item Die Mitgliedstaaten beteiligen sich anteilig an der Finanzierung des Asyl- und Integrationsfonds. Der jährliche Beitrag wird anhand folgender Formel berechnet: (Bruttoinlandsprodukt des Landes in \euro{} x 0,00025).
  
  \item Mitgliedstaaten, die weniger Asylsuchende aufnehmen, als in der Verordnung (EU) 2024/1351 des Europäischen Parlaments und des Rates vom 14.05.2024 festgelegt, müssen pro nicht aufgenommener Person einen zusätzlichen Solidarbeitrag von 20.000 \euro{} in den Asyl- und Integrationsfonds zahlen.

  \item Die Europäische Union gewährt einem Mitgliedstaat, der eine geflüchtete Person aufnimmt, die einem anderen Mitgliedstaat zugeordnet war, einen Zuschuss in Höhe von 10.000 \euro{} pro Person.

  \item Mitgliedstaaten, deren Asyl- und Aufnahmesysteme einem besonderen und unverhältnismäßigen Druck ausgesetzt sind, können auf Antrag mit Mitteln des Asyl- und Integrationsfonds unterstützt werden.

  \item Der Fonds fördert Mitgliedstaaten bei der Gewährung materieller Hilfe, gesundheitlicher Versorgung sowie psychologischer Betreuung für Asylsuchende.

  \item Der Fonds fördert den Aufbau von Kapazitäten zur Umsiedlung von Asylsuchenden aus Mitgliedstaaten, die größere Lasten tragen, in Mitgliedstaaten, die weniger Lasten tragen.
\end{enumerate}

\vspace{-0.2cm}

\subsection*{Artikel 4 – Integration von Schutzsuchenden (EMPL)}

\begin{enumerate}[leftmargin=2em, itemsep=6pt]
  \item Die Mitgliedstaaten ergreifen Maßnahmen zur Förderung der sozialen und wirtschaftlichen Integration von Schutzsuchenden und anerkannten Schutzberechtigten, um deren aktive Teilhabe am gesellschaftlichen Leben sicherzustellen.
  \item Hierzu zählen insbesondere:
  \begin{enumerate}[label=(\alph*), leftmargin=1.8em, itemsep=4pt]
    \item der frühzeitige Zugang zum Arbeitsmarkt sowie die Bereitstellung von Sprachkursen und Qualifizierungsangeboten,

    \item der Zugang zu schulischer, universitärer und beruflicher Bildung,

    \item die Förderung sozialer Inklusion, insbesondere benachteiligter Gruppen,

    \item die Gewährleistung von sozialen Dienstleistungen und der Gesundheitsversorgung,

    \item die Förderung der Mobilität innerhalb der EU für Schutzberechtigte zur Arbeits- oder Ausbildungsaufnahme.
  \end{enumerate}
  \item Die Mitgliedstaaten berichten der Kommission jährlich über ihre Integrationsmaßnahmen.
\end{enumerate}

\subsection*{Artikel 5 – Grenzverfahren (DROI)}

\begin{enumerate}[leftmargin=2em, itemsep=6pt]
  \item Die Asylagentur der Europäischen Union (EUAA) unterstützt die Mitgliedstaaten bei der Identifizierung, Registrierung und Befragung von Asylsuchenden.
  \item Das Asylverfahren an den Außengrenzen der Europäischen Union gliedert sich in zwei aufeinanderfolgende Phasen:
  \begin{enumerate}[label=(\alph*), leftmargin=1.8em, itemsep=4pt]
    \item Vor dem eigentlichen Asylverfahren wird ein Vorprüfungs- und Registrierungsprozess („Screening“) eingeführt:
    \begin{enumerate}[label=(\roman*), leftmargin=1.8em, itemsep=4pt]
      \item Innerhalb einer Frist von mindestens 5 und höchstens 10 Tagen erfolgt eine erste Erfassung der Identität, der Nationalität, des Gesundheitszustands und der Sicherheitsrisiken direkt an den Außengrenzen. Das Ergebnis der Erfassung entscheidet über den weiteren Verfahrensverlauf.
      \item Liegt die Anerkennungsquote des Herkunftslandes unter 20 \% der geflüchteten Personen oder besteht ein relevantes Sicherheitsrisiko, greift das Asylschnellverfahren nach Absatz 2 (b), andernfalls greift das reguläre Asylverfahren.
    \end{enumerate}
    \item Asylschnellverfahren:
    \begin{enumerate}[label=(\roman*), leftmargin=1.8em, itemsep=4pt]
      \item Die Dauer des Asylschnellverfahrens beträgt maximal 12 Wochen und umfasst eine individuelle Prüfung der gestellten Anträge. Jeder antragstellenden Person steht ein kostenloser rechtlicher Beistand zu.
      \item Die Mitgliedstaaten sorgen für die Bereitstellung angemessener Unterbringungsmöglichkeiten, die humanitären Mindeststandards genügen müssen.
      \item In begründeten Einzelfällen kann das Asylschnellverfahren auch in Drittstaaten stattfinden.
    \end{enumerate}
  \end{enumerate}
\end{enumerate}

\subsection*{Artikel 6 – Die Europäische Agentur für die Grenz- und Küstenwache (LIBE)}

\begin{enumerate}[leftmargin=2em, itemsep=6pt]
  \item Die Europäische Agentur für die Grenz- und Küstenwache (FRONTEX) stellt den Mitgliedstaaten auf Antrag operative und technische Unterstützung für die Rückführung abgelehnter Asylbewerberinnen und -bewerber bereit, einschließlich, aber nicht beschränkt auf:
  \begin{enumerate}[label=(\alph*), leftmargin=1.8em, itemsep=4pt]
    \item Ermittlung von Drittstaatsangehörigen ohne Aufenthaltsrecht;
    \item Ausstellung von Rückführungsentscheidungen und Beschaffung notwendiger Reisedokumente;
    \newline
    \item Entsendung von Personal in Drittstaaten zur Unterstützung beim Grenz- und Migrationsmanagement;
    \item Organisation von Rückführungsaktionen von nicht bleibeberechtigten Asylbewerberinnen und -bewerbern in Drittstaaten.
  \end{enumerate}
  \item Die Europäische Agentur für die Grenz- und Küstenwache (FRONTEX) wird bis zum Jahr 2027 mit einer ständigen Reserve von 10.000 FRONTEX-Einsatzkräften ausgestattet, die sicherstellt, dass umfassend auf die Herausforderungen hinsichtlich Grenz- und Migrationsmanagement reagiert werden kann.
\end{enumerate}

\end{document}